%%%%%%%%%%%%%%%%%%%%%%%%%%%%%%%%%%%%%%%%%%%%%%%%%%%%%%%%%%%%%%%%%%%%%
%% Appendix B: Machine Learning and Parametric Modeling Scripts
%% Monograph: Structural Optimization of a 1U CubeSat Chassis with Auxetic Metamaterials
%% Author: Gabriel Marques de Andrade (UESC, 2026)
%%%%%%%%%%%%%%%%%%%%%%%%%%%%%%%%%%%%%%%%%%%%%%%%%%%%%%%%%%%%%%%%%%%%%

\chapter{Machine Learning and Parametric Modeling Scripts}
\label{app:code_listings}

This appendix provides the core Python and PyTorch implementations powering the Physics-Guided Machine Learning (PGML) surrogate pipeline and multi-objective optimization framework described in Chapter~\ref{chap:pgml_optimization}. The implementation adheres to modern scientific computing best practices, including explicit type annotations, modular layer definitions, and rigorous continuum mechanics regularization.

\section{Physics-Guided Deep ResNet PyTorch Architecture}
\label{sec:app_resnet_code}

Listing~\ref{lst:resnet_architecture} details the PyTorch module implementation of the \texttt{ResidualBlock} with Layer Normalization and the complete multi-output \texttt{CubeSatSurrogateResNet} regression architecture.

\begin{lstlisting}[language=Python, caption={PyTorch implementation of the Physics-Guided Deep Residual Network for elastodynamic regression.}, label={lst:resnet_architecture}]
import torch
import torch.nn as nn
from typing import Optional


class ResidualBlock(nn.Module):
    """Dense Residual Block with Layer Normalization and Skip Connection."""

    def __init__(self, hidden_dim: int = 64, dropout_rate: float = 0.05) -> None:
        """Initializes the residual block.

        Args:
            hidden_dim: Number of hidden units in linear layers.
            dropout_rate: Dropout probability for regularization.
        """
        super().__init__()
        self.linear1 = nn.Linear(hidden_dim, hidden_dim)
        self.layer_norm = nn.LayerNorm(hidden_dim)
        self.activation = nn.SiLU()  # Smooth non-linear activation (Swish)
        self.linear2 = nn.Linear(hidden_dim, hidden_dim)
        self.dropout = nn.Dropout(dropout_rate)

    def forward(self, x: torch.Tensor) -> torch.Tensor:
        """Executes forward pass with additive identity shortcut: y = x + F(x)."""
        residual = x
        out = self.linear1(x)
        out = self.layer_norm(out)
        out = self.activation(out)
        out = self.dropout(out)
        out = self.linear2(out)
        return residual + out


class CubeSatSurrogateResNet(nn.Module):
    """Deep Residual MLP for multi-output elastodynamic regression."""

    def __init__(
        self,
        input_dim: int = 5,
        output_dim: int = 4,
        hidden_dim: int = 64,
        num_blocks: int = 3,
    ) -> None:
        """Initializes the surrogate ResNet architecture.

        Args:
            input_dim: Number of input geometric features [theta, t, l, h, rho_rel].
            output_dim: Number of predicted dynamic responses [f1, stress, grms, trans].
            hidden_dim: Width of intermediate dense layers.
            num_blocks: Number of stacked residual blocks.
        """
        super().__init__()
        self.input_dim = input_dim
        self.output_dim = output_dim
        self.hidden_dim = hidden_dim

        # Input linear projection layer
        self.input_layer = nn.Sequential(
            nn.Linear(input_dim, hidden_dim),
            nn.SiLU(),
        )

        # Deep stack of residual blocks with identity shortcuts
        self.res_blocks = nn.ModuleList(
            [ResidualBlock(hidden_dim) for _ in range(num_blocks)]
        )

        # Multi-output regression head
        self.head = nn.Sequential(
            nn.Linear(hidden_dim, hidden_dim // 2),
            nn.SiLU(),
            nn.Linear(hidden_dim // 2, output_dim),
        )

    def forward(self, x: torch.Tensor) -> torch.Tensor:
        """Executes forward inference pass through the deep residual network."""
        out = self.input_layer(x)
        for block in self.res_blocks:
            out = block(out)
        return self.head(out)
\end{lstlisting}

\section{Physics-Guided Loss Function with Gibson-Ashby Penalties}
\label{sec:app_loss_code}

Listing~\ref{lst:physics_guided_loss} presents the custom PyTorch loss module \texttt{PhysicsGuidedLoss}, which combines empirical Mean Squared Error (MSE) with analytical Gibson-Ashby asymptotic power-law penalties and elastodynamic monotonicity constraints.

\begin{lstlisting}[language=Python, caption={Composite Physics-Guided Loss function incorporating Gibson-Ashby and monotonicity penalties.}, label={lst:physics_guided_loss}]
import torch
import torch.nn as nn
from typing import Dict, Optional, Tuple


class PhysicsGuidedLoss(nn.Module):
    """Composite loss function integrating MSE with Gibson-Ashby and Monotonicity penalties."""

    def __init__(
        self,
        lambda_phys: float = 0.15,
        lambda_mono: float = 0.10,
        solid_modulus_gpa: float = 68.0,
    ) -> None:
        """Initializes the physics-guided loss function.

        Args:
            lambda_phys: Weight for the Gibson-Ashby analytical penalty.
            lambda_mono: Weight for the stiffness-density monotonicity constraint.
            solid_modulus_gpa: Young's modulus of bulk AlSi10Mg base alloy.
        """
        super().__init__()
        self.lambda_phys = lambda_phys
        self.lambda_mono = lambda_mono
        self.solid_modulus_gpa = solid_modulus_gpa
        self.mse = nn.MSELoss()

    def forward(
        self,
        y_pred: torch.Tensor,
        y_true: torch.Tensor,
        x_inputs: torch.Tensor,
        target_scaler: Optional[object] = None,
    ) -> Tuple[torch.Tensor, Dict[str, float]]:
        """Computes total physics-guided loss: L_MSE + lambda_phys * L_GA + lambda_mono * L_mono.

        Args:
            y_pred: Predicted target tensor of shape (batch_size, 4).
            y_true: Ground-truth target tensor of shape (batch_size, 4).
            x_inputs: Normalized inputs of shape (batch_size, 5).
            target_scaler: Scaler to un-normalize predictions for physics verification.

        Returns:
            Tuple of (total_loss_tensor, loss_breakdown_dict).
        """
        # 1. Empirical data fidelity loss (Mean Squared Error)
        loss_mse = self.mse(y_pred, y_true)

        # 2. Physics consistency penalty: lower bound (20 Hz) and Gibson-Ashby scaling
        rho_rel = x_inputs[:, 4:5] if x_inputs.shape[1] >= 5 else x_inputs[:, 1:2]
        f1_pred = y_pred[:, 0:1]

        if target_scaler is not None:
            if hasattr(target_scaler, "y_mean") and hasattr(target_scaler, "y_std"):
                mean_f1 = float(target_scaler.y_mean[0])
                std_f1 = float(target_scaler.y_std[0])
            elif hasattr(target_scaler, "mean_y") and hasattr(target_scaler, "std_y"):
                mean_f1 = float(target_scaler.mean_y[0])
                std_f1 = float(target_scaler.std_y[0])
            else:
                mean_f1, std_f1 = None, None

            if mean_f1 is not None and std_f1 is not None:
                f1_physical = f1_pred * std_f1 + mean_f1
                # Penalize predictions violating NASA GEVS excitation minimum (20.0 Hz)
                loss_bound = torch.mean(torch.relu(20.0 - f1_physical) ** 2)

                # Gibson-Ashby asymptotic power-law scaling penalty: f1 ~ sqrt(rho*/rho_s)
                f1_ga = 1350.0 * torch.sqrt(torch.clamp(rho_rel, min=1e-4))
                loss_ga = torch.mean(((f1_physical - f1_ga) / f1_ga) ** 2)
                loss_positivity = loss_bound + 0.10 * loss_ga
            else:
                loss_positivity = torch.mean(torch.relu(-f1_pred) ** 2)
        else:
            loss_positivity = torch.mean(torch.relu(-f1_pred) ** 2)

        # 3. Monotonicity regularization: penalize stiffness drops across density steps
        sorted_indices = torch.argsort(rho_rel.squeeze(-1))
        sorted_f1 = f1_pred[sorted_indices]
        sorted_rho = rho_rel[sorted_indices]
        diff_f1 = sorted_f1[1:] - sorted_f1[:-1]
        diff_rho = sorted_rho[1:] - sorted_rho[:-1]

        # Penalize only where relative density increases significantly (> 0.05)
        significant_step = diff_rho > 0.05
        if significant_step.any():
            severe_drop = torch.relu(-(diff_f1[significant_step] + 0.15))
            loss_mono = (
                torch.mean(severe_drop**2)
                if severe_drop.numel() > 0
                else torch.tensor(0.0, device=y_pred.device)
            )
        else:
            loss_mono = torch.tensor(0.0, device=y_pred.device)

        # Composite total loss
        total_loss = (
            loss_mse
            + self.lambda_phys * loss_positivity
            + self.lambda_mono * loss_mono
        )

        breakdown = {
            "loss_total": float(total_loss.item()),
            "loss_mse": float(loss_mse.item()),
            "loss_positivity": float(loss_positivity.item()),
            "loss_mono": float(loss_mono.item()),
        }
        return total_loss, breakdown
\end{lstlisting}

\section{Feature Normalization and Persistence Scaler}
\label{sec:app_scaler_code}

Listing~\ref{lst:dataset_scaler} provides the data-class implementation of \texttt{SurrogateDatasetScaler}, ensuring strict isolation and bidirectional transformation of training, validation, and test datasets.

\begin{lstlisting}[language=Python, caption={Bidirectional Z-Score dataset normalization scaler and JSON persistence module.}, label={lst:dataset_scaler}]
import json
import logging
from dataclasses import dataclass
from pathlib import Path
from typing import List
import numpy as np

logger = logging.getLogger(__name__)


@dataclass
class SurrogateDatasetScaler:
    """Standard Z-Score normalization scaler for neural inputs and elastodynamic targets."""

    x_mean: np.ndarray
    x_std: np.ndarray
    y_mean: np.ndarray
    y_std: np.ndarray
    input_names: List[str]
    target_names: List[str]

    def transform_x(self, x: np.ndarray) -> np.ndarray:
        """Transforms raw geometric inputs to normalized Z-scores."""
        return (x - self.x_mean) / np.maximum(self.x_std, 1e-8)

    def inverse_transform_x(self, z_x: np.ndarray) -> np.ndarray:
        """Inverse-transforms Z-scores back to physical geometric dimensions."""
        return z_x * self.x_std + self.x_mean

    def transform_y(self, y: np.ndarray) -> np.ndarray:
        """Transforms raw elastodynamic targets to normalized Z-scores."""
        return (y - self.y_mean) / np.maximum(self.y_std, 1e-8)

    def inverse_transform_y(self, z_y: np.ndarray) -> np.ndarray:
        """Inverse-transforms Z-scores back to physical engineering units."""
        return z_y * self.y_std + self.y_mean

    def save_json(self, filepath: str | Path) -> Path:
        """Saves scaler statistical parameters to a JSON file."""
        target_path = Path(filepath)
        target_path.parent.mkdir(parents=True, exist_ok=True)
        data = {
            "x_mean": self.x_mean.tolist(),
            "x_std": self.x_std.tolist(),
            "y_mean": self.y_mean.tolist(),
            "y_std": self.y_std.tolist(),
            "input_names": self.input_names,
            "target_names": self.target_names,
        }
        with open(target_path, "w", encoding="utf-8") as f:
            json.dump(data, f, indent=2)
        logger.info("Saved scaler parameters to %s", target_path)
        return target_path

    @classmethod
    def load_json(cls, filepath: str | Path) -> "SurrogateDatasetScaler":
        """Loads scaler statistical parameters from a JSON configuration file."""
        with open(filepath, "r", encoding="utf-8") as f:
            data = json.load(f)
        return cls(
            x_mean=np.array(data["x_mean"], dtype=np.float32),
            x_std=np.array(data["x_std"], dtype=np.float32),
            y_mean=np.array(data["y_mean"], dtype=np.float32),
            y_std=np.array(data["y_std"], dtype=np.float32),
            input_names=data["input_names"],
            target_names=data["target_names"],
        )
\end{lstlisting}
